% =====================================================================
%  CARNET D'ENTRAÎNEMENT — Fiche 12 (Chimie organique)
%  THÈME 4 — Masse molaire d'une molécule organique
%  NIVEAU MOYEN — ÉNONCÉS
% =====================================================================
\documentclass[11pt,a4paper]{article}
\usepackage[utf8]{inputenc}
\usepackage[T1]{fontenc}
\usepackage{lmodern}
\usepackage[french]{babel}
\usepackage{amsmath,amssymb,mathtools}
\providecommand{\degree}{^\circ}
\usepackage{icomma}
\usepackage{siunitx}
\sisetup{output-decimal-marker={,},per-mode=symbol}
\usepackage[version=4]{mhchem}
\usepackage{chemfig}
\setchemfig{atom sep=2em,bond length=0.9em,cram width=2.5pt}
\usepackage{xcolor}
\usepackage{array}
\usepackage{booktabs}
\usepackage{tabularx}
\usepackage[most]{tcolorbox}
\usepackage{enumitem}
\usepackage{tikz}
\usetikzlibrary{arrows.meta,positioning,calc}
\usepackage{fancyhdr}
\usepackage{titlesec}
\usepackage[hidelinks]{hyperref}
\usepackage[a4paper,margin=2.1cm,headheight=26pt,headsep=9pt]{geometry}

\definecolor{primary}{RGB}{150,98,162}
\definecolor{primarydark}{RGB}{92,52,110}
\definecolor{excol}{RGB}{0,135,90}
\definecolor{attncol}{RGB}{200,40,55}
\definecolor{methcol}{RGB}{90,90,100}

\newcommand{\runtitle}{Thème 4 — Moyen — Énoncés}
\pagestyle{fancy}\fancyhf{}
\fancyhead[L]{\small\textcolor{primarydark}{\textbf{Fiche 12} — Chimie organique}}
\fancyhead[R]{\small\textcolor{primarydark}{\runtitle}}
\fancyfoot[C]{\small\thepage}
\renewcommand{\headrulewidth}{0.8pt}
\renewcommand{\headrule}{\hbox to\headwidth{\color{primary}\leaders\hrule height \headrulewidth\hfill}}
\setlength{\parindent}{0pt}\setlength{\parskip}{3pt}

\tcbset{boxbase/.style={enhanced,breakable,boxrule=0.9pt,arc=2.5pt,
   left=3mm,right=3mm,top=2mm,bottom=2mm,fonttitle=\bfseries,coltitle=white,
   toptitle=0.5mm,bottomtitle=0.5mm}}
\newtcolorbox[auto counter]{exo}[1]{boxbase,colback=primary!4,colframe=primary,
   title={\textsc{Exercice}~\thetcbcounter\ \textmd{--- \itshape #1}}}
\newtcolorbox{consigne}{boxbase,colback=methcol!8,colframe=methcol,sharp corners,
   title={\textsc{Consignes \& données}}}

\newcommand{\banner}[2]{%
\begin{tcolorbox}[enhanced,colback=primary,colframe=primary,arc=5pt,boxrule=0pt,
   left=5mm,right=5mm,top=2.5mm,bottom=2.5mm,coltext=white]
{\large\bfseries #1}\\[1pt]{\small #2}
\end{tcolorbox}\vspace{2pt}}

\begin{document}

\banner{Thème 4 — Masse molaire d'une molécule organique}{Niveau \textbf{Moyen} — Énoncés \quad$\vert$\quad Fiche 12, Chimie organique}

\begin{consigne}
Exercices \textbf{ouverts}. \textbf{Données} : $M(\ce{H})=1$, $M(\ce{C})=12$, $M(\ce{N})=14$, $M(\ce{O})=16$ (en $\si{\g\per\mol}$) ; $N_A = \qty{6,02e23}{\per\mol}$. Convertis masses et volumes avant tout calcul.
\end{consigne}

% ---------------------------------------------------------------------
\begin{exo}{masse molaire du paracétamol}
Le paracétamol a pour formule brute $\ce{C8H9NO2}$. Calcule sa masse molaire $M$ (contribution par contribution).
\end{exo}

% ---------------------------------------------------------------------
\begin{exo}{nombre de molécules (aspirine)}
L'aspirine ($\ce{C9H8O4}$) a pour masse molaire $M=\qty{180}{\g\per\mol}$. Combien de \textbf{molécules} d'aspirine contient un comprimé de $\qty{360}{\mg}$ ?
\end{exo}

% ---------------------------------------------------------------------
\begin{exo}{concentration d'une solution (paracétamol)}
On dissout $\qty{755}{\mg}$ de paracétamol ($M=\qty{151}{\g\per\mol}$) dans $\qty{200}{\mL}$ de solution. Calcule la \textbf{concentration molaire} $C$ du paracétamol.
\end{exo}

% ---------------------------------------------------------------------
\begin{exo}{quelle masse peser ? (mannitol)}
Le mannitol (un édulcorant, un hexol) a pour formule brute $\ce{C6H14O6}$.
\begin{enumerate}[label=\textbf{\alph*)},leftmargin=2em,itemsep=1pt]
  \item Calcule sa masse molaire $M$.
  \item Quelle \textbf{masse} de mannitol faut-il peser pour disposer de $\qty{1,20e21}{}$ molécules ?
\end{enumerate}
\end{exo}

% ---------------------------------------------------------------------
\begin{exo}{le piège de l'unité}
Une molécule a pour formule brute $\ce{C6H14N2}$.
\begin{enumerate}[label=\textbf{\alph*)},leftmargin=2em,itemsep=1pt]
  \item Calcule sa \textbf{masse moléculaire relative} $M_r$.
  \item Un candidat répond «\,$M_r = \qty{114}{\g\per\mol}$\,». Est-ce correct ? Justifie et donne la réponse juste.
\end{enumerate}
\end{exo}

\end{document}
